\documentclass[preview]{standalone}

\usepackage{amsmath}
\usepackage{amssymb}
\usepackage{stellar}
\usepackage{definitions}

\begin{document}

\id{permutation-conjugacy}
\genpage

\section{Conjugacy in symmetric groups}

\begin{snippettheorem}{symmetric-group-conjugacy-cycle-type}{Conjugacy and cycle type}
    Two permutations in \(\permgrp_n\) are conjugate \ifandonlyif they have
    the same cycle type.
\end{snippettheorem}

\begin{snippetproof}{symmetric-group-conjugacy-cycle-type-proof}{symmetric-group-conjugacy-cycle-type}{Conjugacy and cycle type}
    Conjugating a cycle relabels its entries:
    \[
        g^{-1}(a_1\ a_2\ \cdots\ a_r)g
        =(g^{-1}(a_1)\ g^{-1}(a_2)\ \cdots\ g^{-1}(a_r)).
    \]
    Thus conjugation preserves the lengths and multiplicities of the disjoint
    cycles. Conversely, if \(\sigma\) and \(\tau\) have the same cycle type,
    pair their cycles of equal length and choose a bijection \(g\) carrying
    corresponding entries in cyclic order. The displayed identity then gives
    \(g^{-1}\sigma g=\tau\).
\end{snippetproof}

\begin{snippetcorollary}{symmetric-group-conjugacy-class-size}{Conjugacy class size in a symmetric group}
    Let \(\sigma\in\permgrp_n\) have \(m_r\) cycles of length \(r\), counting
    fixed points as cycles of length \(1\). Then
    \[
        |\groupcentralizer_{\permgrp_n}(\sigma)|
        =\prod_{r\geq1}r^{m_r}m_r!,
        \qquad
        |\sigma^{\permgrp_n}|
        =\frac{n!}{\prod_{r\geq1}r^{m_r}m_r!}.
    \]
\end{snippetcorollary}

\begin{snippetproof}{symmetric-group-conjugacy-class-size-proof}{symmetric-group-conjugacy-class-size}{Conjugacy class size in a symmetric group}
    A permutation commuting with \(\sigma\) may independently rotate every
    \(r\)-cycle in \(r\) ways and permute the \(m_r\) cycles of the same
    length in \(m_r!\) ways. These choices describe the entire centralizer.
    The class-size formula follows from orbit--stabilizer.
\end{snippetproof}

\section{Conjugacy in alternating groups}

\begin{snippettheorem}{alternating-group-conjugacy-class-splitting}{Splitting of conjugacy classes in \(A_n\)}
    Let \(\sigma\in A_n\). Its conjugacy class in \(\permgrp_n\) is either a
    single conjugacy class in \(A_n\), or the disjoint union of two
    \(A_n\)-conjugacy classes of equal size. It splits into two classes if
    and only if the cycle lengths of \(\sigma\), including \(1\)-cycles, are
    odd and pairwise distinct.
\end{snippettheorem}

\begin{snippetproof}{alternating-group-conjugacy-class-splitting-proof}{alternating-group-conjugacy-class-splitting}{Splitting of conjugacy classes in \(A_n\)}
    Since \(A_n\) has index \(2\) in \(\permgrp_n\), the
    \(\permgrp_n\)-orbit of \(\sigma\) consists of one or two \(A_n\)-orbits.
    The two possible orbits coincide exactly when the centralizer
    \(\groupcentralizer_{\permgrp_n}(\sigma)\) contains an odd permutation:
    such a centralizing element changes an odd conjugator into an even one
    without changing the conjugate. If no such element exists, the even and
    odd conjugators give two orbits of equal size.

    It remains to characterize when the centralizer contains an odd element.
    An even-length cycle is itself odd and centralizes \(\sigma\). If two
    cycles have the same odd length, the permutation exchanging those two
    cycles is odd and centralizes \(\sigma\). Thus splitting requires odd,
    pairwise distinct cycle lengths. Conversely, under that condition every
    centralizing permutation must preserve each cycle and act on it by a
    rotation. Powers of odd-length cycles are even, so every element of the
    centralizer is even. This proves the criterion.
\end{snippetproof}

\end{document}
